\documentclass{exam}
\usepackage{amsmath, amssymb}

\pagestyle{headandfoot}
\firstpageheadrule
\runningheadrule
\firstpageheader{Prof. Tahvildar-Zadeh \\ Linear Algebra}{Homework 11}{Aadhithya Saravanan}
\runningheader{Linear Algebra \\ Homework 11}{}{Saravanan}
\firstpagefooter{}{}{}
\runningfooter{}{\thepage}{}

\printanswers

\begin{document}

\underline{Exercise 6.1.2}
\newline
\newline
\[\langle x,y\rangle=2(2+i)+(i+1)(2)+i(1-2i)\]
\[=4+2i+2i+2+i+2\]
\[=8+5i\]
\newline
\[\|x\|=\sqrt{\langle x,x\rangle}\]
\[=\sqrt{2^2+(1+i)(1-i)+i(-i)}\]
\[=\sqrt{4+2+1}\]
\[=\sqrt{7}\]
\newline
\[\|y\|=\sqrt{\langle y,y\rangle}\]
\[=\sqrt{(2-i)(2+i)+2^2+(1+2i)(1-2i)}\]
\[=\sqrt{5+4+5}\]
\[=\sqrt{14}\]
\newline
\[\|x+y\|=\sqrt{\langle x+y,x+y\rangle}\]
\[=\sqrt{(4-i)(4+i)+(3+i)(3-i)+(1+3i)(1-3i)}\]
\[=\sqrt{17+10+10}\]
\[=\sqrt{37}\]
\newline
Cauchy-Schwarz Inequality:
\[|\langle x,y\rangle|\le\|x\|\|y\|\]
\[|8+5i|\le\sqrt{7}\sqrt{14}\]
\[\sqrt{8^2+5^2}\le\sqrt{98}\]
\[\sqrt{64+25}\le\sqrt{98}\]
\[\sqrt{89}\le\sqrt{98}\]
\newline
Triangle Inequality:
\[\|x+y\|\le\|x\|+\|y\|\]
\[\sqrt{37}\le\sqrt{7}+\sqrt{14}\]
Squaring both sides:
\[37\le7+14+2\sqrt{98}\]
Subtracting 21 from both sides and dividing by 2:
\[8\le\sqrt{98}\]
\[8=\sqrt{64}\le\sqrt{98}\]
\newline

\underline{Exercise 6.1.3}
\newline
\newline
\[\langle f,g\rangle=\int_{0}^{1}te^tdt\]
\[=\left[te^t-\int e^tdt\right]_0^1\]
\[=\left[te^t-e^t\right]_0^1\]
\[=e-e-0+1\]
\[=1\]
\newline
\[\|f\|=\sqrt{\langle f,f\rangle}\]
\[=\sqrt{\int_{0}^{1}t^2dt}\]
\[=\sqrt{\frac{1}{3}}\]
\[=\frac{\sqrt{3}}{3}\]
\newline
\[\|g\|=\sqrt{\langle g,g\rangle}\]
\[=\sqrt{\int_{0}^{1}e^{2t}dt}\]
\[=\sqrt{\frac{1}{2}(e^2-1)}\]
\[=\frac{\sqrt{2e^2-2}}{2}\]
\newline
\[\|f+g\|=\sqrt{\langle f+g,f+g\rangle}\]
\[=\sqrt{\int_{0}^{1}(t+e^t)(t+e^t)dt}\]
\[=\sqrt{\int_{0}^{1}t^2+2te^t+e^{2t}dt}\]
\[=\sqrt{\int_{0}^{1}t^2dt+2\int_{0}^{1}te^tdt+\int_{0}^{1}e^{2t}dt}\]
\[=\sqrt{\frac{1}{3}+2+\frac{1}{2}(e^2-1)}\]
\[=\sqrt{\frac{1}{2}e^2+\frac{11}{6}}\]
\[=\frac{\sqrt{18e^2+66}}{6}\]
\newline
Cauchy-Schwarz Inequality:
\[|\langle f,g\rangle|\le\|f\|\|g\|\]
\[|1|\le\frac{\sqrt{3}}{3}\frac{\sqrt{2e^2-2}}{2}\]
\[1\le\frac{\sqrt{6e^2-6}}{6}\]
Multiplying both sides by 6 and squaring:
\[36\le6e^2-6\]
Adding 6 to both sides and dividing by 6:
\[7\le e^2\]
Using the first few terms of the infinite series expansion for $e$:
\[e>1+1+\frac{1}{2}+\frac{1}{6}=\frac{8}{3}\]
So,
\[7<\frac{64}{9}=\left(\frac{8}{3}\right)^2<e^2\]
\newline
Triangle Inequality:
\[\|f+g\|\le\|f\|+\|g\|\]
\[\frac{\sqrt{18e^2+66}}{6}\le\frac{\sqrt{3}}{3}+\frac{\sqrt{2e^2-2}}{2}\]
Multiplying both sides by 6 and squaring:
\[18e^2+66\le12+9(2e^2-2)+2(2\sqrt{3})(3\sqrt{2e^2-2})\]
\[18e^2+66\le18e^2-6+12\sqrt{6e^2-6}\]
Subtracting $18e^2$ and adding 6 to each side:
\[72\le12\sqrt{6e^2-6}\]
Dividing both sides by 12 and squaring:
\[36\le6e^2-6\]
Adding 6 to each side and dividing by 6:
\[7\le e^2\]
Using the first few terms of the infinite series expansion for $e$:
\[e>1+1+\frac{1}{2}+\frac{1}{6}=\frac{8}{3}\]
So,
\[7<\frac{64}{9}=\left(\frac{8}{3}\right)^2<e^2\]
\newline

\underline{Exercise 6.1.9}
\newline
\begin{parts}
    \part Since the vector space is finite-dimensional, let its dimension be $n$. Then, $\beta$ has $n$ vectors. Let $\beta=\{z_1,z_2,\dots,z_n\}$. We want to prove that $x=0$ given that $\langle x,z_i\rangle=0$ for all $z_i\in\beta$. Let $z$ be in the vector space. By definition of a basis, we can write $z$ as $\sum_{i=1}^{n}a_iz_i$, where each $a_i$ is a scalar. Then, by Theorem 6.1.a and 6.1.b, we have $\langle x,z\rangle=\langle x,\sum_{i=1}^{n}a_iz_i\rangle=\sum_{i=1}^{n}\overline{a_i}\langle x,z_i\rangle=\sum_{i=1}^{n}\overline{a_i}(0)=0$ for all $z$. Let $z=x$. Then $\langle x,x\rangle=0$. By Theorem 6.1.d, $x=0$.
    \part Since the vector space is finite-dimensional, let its dimension be $n$. Then, $\beta$ has $n$ vectors. Let $\beta=\{z_1,z_2,\dots,z_n\}$. We want to prove that $x=y$ given that $\langle x,z_i\rangle=\langle y,z_i\rangle$ for all $z_i\in\beta$. So, $\langle x,z_i\rangle-\langle y,z_i\rangle=0$. By definition of an inner product, we have $\langle x,z_i\rangle-\langle y,z_i\rangle=\langle x-y,z_i\rangle$. Then, $\langle x-y,z_i\rangle=0$. Let $z$ be in the vector space. By definition of a basis, we can write $z$ as $\sum_{i=1}^{n}a_iz_i$, where each $a_i$ is a scalar. Then, by Theorem 6.1.a and 6.1.b, we have $\langle x-y,z\rangle=\langle x-y,\sum_{i=1}^{n}a_iz_i\rangle=\sum_{i=1}^{n}\overline{a_i}\langle x-y,z_i\rangle=\sum_{i=1}^{n}\overline{a_i}(0)=0$ for all $z$. Choose $z=x-y$. Then, $\langle x-y,x-y\rangle=0$. By Theorem 6.1.d, $x-y=0$, so $x=y$.
    \newline
\end{parts}

\underline{Exercise 6.1.11}
\newline
\newline
\[\|x+y\|^2+\|x-y\|^2\]
\[=(\sqrt{\langle x+y,x+y\rangle})^2+(\sqrt{\langle x-y,x-y\rangle})^2\]
\[=\langle x+y,x+y\rangle+\langle x-y,x-y\rangle\]
By Theorem 6.1.a and Theorem 6.1.b,
\[=\langle x+y,x\rangle+\langle x+y,y\rangle+\langle x-y,x\rangle-\langle x-y,y\rangle\]
By definition of an inner product,
\[=\langle x,x\rangle+\langle y,x\rangle+\langle x,y\rangle+\langle y,y\rangle+\langle x,x\rangle-\langle y,x\rangle-\langle x,y\rangle+\langle y,y\rangle\]
\[=2\langle x,x\rangle+2\langle y,y\rangle\]
\[=2(\sqrt{\langle x,x\rangle})^2+2(\sqrt{\langle y,y\rangle})^2\]
\[=2\|x\|^2+2\|y\|^2\]
In $\mathbb{R}^2$, let $x$ and $y$ be vectors representing the adjacent sides of a parallelogram. Then, $x+y$ and $x-y$ are the diagonals of the parallelogram. The parallelogram law then says that the sum of the squares of the lengths of the diagonals is equal to the sum of the squares of the lengths of all four sides.
\newline

\underline{Exercise 6.2.2.a}
\newline
\newline
Let $S=\{(1,0,1),(0,1,1),(1,3,3)\}=\{w_1,w_2,w_3\}$. We are looking for an orthogonal basis consisting of $v_1,v_2,v_3$. Let $v_1=w_1$. Then,
\[v_2=w_2-\frac{\langle w_2,v_1\rangle}{\|v_1\|^2}v_1\]
\[=(0,1,1)-\frac{1}{2}(1,0,1)\]
\[=\left(-\frac{1}{2},1,\frac{1}{2}\right)\]
Now, we calculate $v_3$ using $w_3$, $v_1$, and $v_2$.
\[v_3=w_3-\frac{\langle w_3,v_1\rangle}{\|v_1\|^2}v_1-\frac{\langle w_3,v_2\rangle}{\|v_2\|^2}v_2\]
\[=(1,3,3)-\frac{4}{2}(1,0,1)-\frac{4}{\frac{3}{2}}\left(-\frac{1}{2},1,\frac{1}{2}\right)\]
\[=\left(\frac{1}{3},\frac{1}{3},-\frac{1}{3}\right)\]
So, $\alpha=\{v_1,v_2,v_3\}=\{(1,0,1),\left(-\frac{1}{2},1,\frac{1}{2}\right),\left(\frac{1}{3},\frac{1}{3},-\frac{1}{3}\right)\}$ is an orthogonal basis for $\operatorname{span}(S)$. We can normalize each of these vectors to obtain an orthonormal basis $\beta$.
\[\beta_1=\frac{1}{\|v_1\|}v_1,\beta_2=\frac{1}{\|v_2\|}v_2,\beta_3=\frac{1}{\|v_3\|}v_3\]
\[\beta_1=\frac{1}{\sqrt{\langle v_1,v_1\rangle}}v_1,\beta_2=\frac{1}{\sqrt{\langle v_2,v_2\rangle}}v_2,\beta_3=\frac{1}{\sqrt{\langle v_3,v_3\rangle}}v_3\]
\[\beta_1=\frac{1}{\sqrt{2}}(1,0,1),\beta_2=\frac{1}{\sqrt{\frac{3}{2}}}\left(-\frac{1}{2},1,\frac{1}{2}\right),\beta_3=\frac{1}{\sqrt{\frac{1}{3}}}\left(\frac{1}{3},\frac{1}{3},-\frac{1}{3}\right)\]
\[\beta_1=\left(\frac{\sqrt{2}}{2},0,\frac{\sqrt{2}}{2}\right),\beta_2=\left(-\frac{\sqrt{6}}{6},\frac{\sqrt{6}}{3},\frac{\sqrt{6}}{6}\right),\beta_3=\left(\frac{\sqrt{3}}{3},\frac{\sqrt{3}}{3},-\frac{\sqrt{3}}{3}\right)\]
So, $\beta=\{\beta_1,\beta_2,\beta_3\}$ is an orthonormal basis for $\operatorname{span}(S)$. The Fourier coefficients of $x=(1,1,2)$ relative to $\beta$ are the following:
\[(\langle x,\beta_1\rangle,\langle x,\beta_2\rangle,\langle x,\beta_3\rangle)\]
\[=\left(\frac{3\sqrt{2}}{2},\frac{\sqrt{6}}{2},0\right)\]
To verify by Theorem 6.5, we compute the following:
\[\sum_{i=1}^{3}\langle x,\beta_i\rangle\beta_i\]
\[=\frac{3\sqrt{2}}{2}\left(\frac{\sqrt{2}}{2},0,\frac{\sqrt{2}}{2}\right)+\frac{\sqrt{6}}{2}\left(-\frac{\sqrt{6}}{6},\frac{\sqrt{6}}{3},\frac{\sqrt{6}}{6}\right)+0\left(\frac{\sqrt{3}}{3},\frac{\sqrt{3}}{3},-\frac{\sqrt{3}}{3}\right)\]
\[=\left(\frac{3}{2},0,\frac{3}{2}\right)+\left(-\frac{1}{2},1,\frac{1}{2}\right)+(0,0,0)\]
\[=(1,1,2)\]
\[=x\]
Since this is equal to the original vector $x$, we have verified our calculation.
\newline

\underline{Exercise 6.2.2.c}
\newline
\newline
Let $S=\{1,x,x^2\}=\{w_1,w_2,w_3\}$. We are looking for an orthogonal basis consisting of $v_1,v_2,v_3$. Let $v_1=w_1$. Then,
\[v_2=w_2-\frac{\langle w_2,v_1\rangle}{\|v_1\|^2}v_1\]
\[=x-\frac{\frac{1}{2}}{1}1\]
\[=x-\frac{1}{2}\]
Now, we calculate $v_3$ using $w_3$, $v_1$, and $v_2$.
\[v_3=w_3-\frac{\langle w_3,v_1\rangle}{\|v_1\|^2}v_1-\frac{\langle w_3,v_2\rangle}{\|v_2\|^2}v_2\]
\[=x^2-\frac{\frac{1}{3}}{1}1-\frac{\frac{1}{12}}{\frac{1}{12}}\left(x-\frac{1}{2}\right)\]
\[=x^2-x+\frac{1}{6}\]
So, $\alpha=\{v_1,v_2,v_3\}=\{1,x-\frac{1}{2},x^2-x+\frac{1}{6}\}$ is an orthogonal basis for $\operatorname{span}(S)$. We can normalize each of these vectors to obtain an orthonormal basis $\beta$.
\[\beta_1=\frac{1}{\|v_1\|}v_1,\beta_2=\frac{1}{\|v_2\|}v_2,\beta_3=\frac{1}{\|v_3\|}v_3\]
\[\beta_1=\frac{1}{\sqrt{\langle v_1,v_1\rangle}}v_1,\beta_2=\frac{1}{\sqrt{\langle v_2,v_2\rangle}}v_2,\beta_3=\frac{1}{\sqrt{\langle v_3,v_3\rangle}}v_3\]
\[\beta_1=\frac{1}{1}(1),\beta_2=\frac{1}{\sqrt{\frac{1}{12}}}\left(x-\frac{1}{2}\right),\beta_3=\frac{1}{\sqrt{\frac{1}{180}}}\left(x^2-x+\frac{1}{6}\right)\]
\[\beta_1=1,\beta_2=2\sqrt{3}x-\sqrt{3},\beta_3=6\sqrt{5}x^2-6\sqrt{5}x+\sqrt{5}\]
So, $\beta=\{\beta_1,\beta_2,\beta_3\}$ is an orthonormal basis for $\operatorname{span}(S)$. The Fourier coefficients of $h(x)=1+x$ relative to $\beta$ are the following:
\[(\langle h,\beta_1\rangle,\langle h,\beta_2\rangle,\langle h,\beta_3\rangle)\]
\[=\left(\frac{3}{2},\frac{\sqrt{3}}{6},0\right)\]
To verify by Theorem 6.5, we compute the following:
\[\sum_{i=1}^{3}\langle h,\beta_i\rangle\beta_i\]
\[=\frac{3}{2}(1)+\frac{\sqrt{3}}{6}\left(2\sqrt{3}x-\sqrt{3}\right)+0\left(6\sqrt{5}x^2-6\sqrt{5}x+\sqrt{5}\right)\]
\[=\frac{3}{2}+\left(x-\frac{1}{2}\right)+0\]
\[=1+x\]
Since this is equal to the original vector $h(x)$, we have verified our calculation.
\newline

\end{document}